\section{Method}
\label{sec:method}
\subsection{Two-Stage Training}
The pipeline has two stages, both using ResNet-18 \cite{he2016resnet}. The first
stage fits a 241-way classifier on \cnPreTrainN{} images for \cnPreEpochs{}
epochs starting from the ImageNet weights of our Western baseline, so that the
trunk has already seen Chinese food before it sees the twenty dishes. The second
stage replaces the classifier with a 21-way head and fine-tunes for
\cnEpochs{} epochs from the first stage's best checkpoint, keeping the last
residual stage and the head trainable and freezing the rest. Both stages share
one recipe: AdamW with weight decay $5\times10^{-4}$, head learning rate
$10^{-3}$, backbone learning rate $2\times10^{-4}$, cosine annealing over the
epoch budget, batch size 64, label smoothing 0.1
\cite{szegedy2016inception}, mixed precision, and input resolution
$224\times224$. Training uses a random resized crop to $224\times224$ with scale
range $[0.55,1.0]$, a horizontal flip, and a photometric jitter with brightness,
contrast and saturation in $[0.65,1.35]$ and hue in $[-0.06,0.06]$. Checkpoints
are selected by validation accuracy and evaluated once on the test split.

The two stages are deliberately identical apart from the label space and the
initialisation, which is what makes the pre-training comparison in
Section \ref{sec:pretrain} a controlled one: the same images, the same schedule
and the same head capacity, differing only in whether the trunk arrived with
Chinese dish features or with ImageNet features.

\subsection{Six-View Test-Time Fusion}
A tray photograph has no canonical crop. The plate may fill the frame, or the
frame may include the whole tray with rice, soup and a table edge. Instead of
committing to one preprocessing convention, we average the posteriors of six
deterministic views: a $256$-pixel short-side resize with a $224$ centre crop, a
direct $224\times224$ resize, and a $320$-pixel short-side resize with a $224$
centre crop, each also in horizontal mirror. The six posteriors are averaged in
probability space, and the result is divided by a temperature fitted on the
validation split. We report both the single-view protocol used for model
selection and the fused protocol, because the difference between them is a
deployment choice rather than a property of the weights.

\subsection{Open-Set Gate}
\label{sec:gate}
A softmax head always answers, so the deployed system decides separately whether
to answer at all. Three signals are combined, all computed from the same forward
passes:

\textbf{Domain probe.} A logistic regression on the pooled 512-dimensional
features of the three unmirrored views, trained on canteen training images
against a mixture of non-canteen photographs, outputs the probability that the
input belongs to the canteen domain at all. Because it reads pooled features
rather than logits, it is cheap and it does not need the twenty-one classes to be
correct.

\textbf{Chinese-dish veto.} A second, wider head covers 116 Chinese dishes
including our twenty. If that head, with the canteen classes masked out,
confidently names a dish that the canteen head did not, the input is more likely
to be one of the other ninety-six dishes, and the answer is withheld.

\textbf{Catch-all posterior.} The twenty-first class is a trained outcome, so a
high posterior on it is itself evidence that the plate holds something outside
the menu. A separate food-versus-non-food gate and a packaging detector remove
photographs of menus, wrappers and empty trays before any dish is named.

A photograph is answered only if the fused posterior passes a threshold on the
twenty-dish mass, the domain probe is above its threshold, the veto did not fire,
the catch-all posterior is low, and both auxiliary gates agree. Every threshold
is fitted on the canteen validation split, never on the test split, and the
coverage that this costs on supported inputs is reported together with the
rejections it buys in Section \ref{sec:openset}.

\subsection{Energy Interval}
\label{sec:calorie}
Let $k_c$ be the energy density of class $c$ in kcal per 100\,g and $m_c$ its
typical portion mass in grams. The point estimate is
\begin{equation}
\hat{E}_c = \frac{k_c\, m_c}{100}.
\end{equation}
Because the portion is unknown, we report
\begin{equation}
\hat{E}_c\left[1-s_c,\; 1+s_c\right], \quad
s_c = 0.15+0.10\cdot\mathbf{1}[p_{\max}<0.55],
\label{eq:interval}
\end{equation}
where $p_{\max}$ is the posterior of the top-1 class. The first term
acknowledges recipe and portion variability, the second widens the interval when
the classifier itself is uncertain, so that a wrong but confident guess is not
presented as a precise number. The interface offers a small, regular or large
portion, which scales $m_c$ by $0.7$, $1$ or $1.35$, and accepts a measured mass
when the user has one, in which case only the density prior is used.

\subsection{Metrics and Complexity}
We report top-1 and top-5 accuracy, macro-F1 over the twenty-one classes, and
expected calibration error with fifteen equal-width confidence bins, together
with parameter count, multiply-accumulate operations at $224\times224$ measured
with THOP, and single-image latency on the GPU and on the CPU after warm-up. For
the gate we report acceptance rate per evaluation set rather than a single
accuracy, because a gate that accepts everything and a gate that accepts nothing
have the same accuracy on a set of unseen dishes.
